\chapter*{Introduction générale}
\markboth{Introduction générale}{Introduction générale}

\section*{Contexte}

À l'ère de la transformation numérique, la majorité des services
quotidiens — communication, gestion documentaire, commerce,
administration — repose sur des applications web hébergées dans le
\emph{cloud}. Cette omniprésence du Web s'est accompagnée d'une
explosion des incidents de sécurité : selon le rapport
\textit{Verizon Data Breach Investigations Report 2023}, plus de 80\,\%
des compromissions enregistrées exploitent des vulnérabilités
applicatives ou des erreurs de configuration. Concevoir une application
web ne se résume donc plus à coder un service fonctionnel : il faut
également maîtriser son hébergement, son réseau, son cycle de
déploiement et sa surface d'attaque.

Le mouvement \emph{DevOps} et l'adoption généralisée des \textbf{conteneurs}
ont profondément changé la manière de construire et de livrer ces
applications. Docker offre une portabilité quasi parfaite entre les
environnements de développement et de production, tandis que les
fournisseurs de \emph{cloud} comme \textbf{Amazon Web Services} (AWS)
permettent de provisionner en quelques minutes des infrastructures
auparavant réservées aux grands centres de données. La maîtrise de
ces outils est désormais une compétence attendue de tout développeur
ou administrateur système.

\section*{Problématique}

Si les briques technologiques sont accessibles, leur \emph{intégration
sécurisée} reste un exercice difficile : combien de bases de données
exposées sur Internet, combien de certificats expirés, combien
d'applications vulnérables aux injections SQL ou au \emph{cross-site
scripting} ?

La problématique de ce projet peut être formulée ainsi :

\begin{quote}
\itshape Comment concevoir, développer et déployer une application
web fonctionnelle dans le \emph{cloud}, en appliquant de manière
explicite et vérifiable les bonnes pratiques de cybersécurité, depuis
la couche réseau jusqu'au code applicatif ?
\end{quote}

\section*{Objectifs}

Le projet \textbf{WebCyber} apporte une réponse pédagogique et
pragmatique à cette problématique en poursuivant les objectifs
suivants :

\begin{enumerate}
  \item Concevoir une application web simple mais réaliste : une
        application personnelle de prise de notes (CRUD complet,
        authentification, isolation des données par utilisateur).
  \item Conteneuriser l'ensemble des composants (Nginx, application
        Flask, base PostgreSQL) avec Docker et Docker Compose.
  \item Déployer la solution sur une infrastructure AWS minimaliste
        mais représentative (EC2, Elastic IP, Security Group, Route\,53).
  \item Mettre en œuvre une chaîne de sécurité \textbf{en profondeur}
        couvrant le réseau, le système, les conteneurs, le proxy, le
        code applicatif et la base de données.
  \item Activer HTTPS via un certificat Let's Encrypt, avec
        renouvellement automatisé.
  \item Évaluer la posture de sécurité du déploiement à l'aide
        d'outils standard (Nmap, Nikto, testssl.sh) et discuter les
        résultats.
  \item Rédiger une documentation technique exploitable permettant à
        un autre étudiant ou ingénieur de reproduire la solution.
\end{enumerate}

\section*{Plan du rapport}

Ce mémoire s'articule autour de cinq chapitres :

\begin{itemize}
  \item Le \textbf{chapitre 1} pose le \textit{contexte général} et
        l'\textit{état de l'art} : \emph{cloud computing},
        conteneurisation, architecture des applications web modernes
        et principales menaces de sécurité (OWASP Top 10).
  \item Le \textbf{chapitre 2} formalise l'\textit{analyse et la
        spécification des besoins} : acteurs, cas d'utilisation,
        besoins fonctionnels et non fonctionnels.
  \item Le \textbf{chapitre 3} présente la \textit{conception} de la
        solution : architecture globale, infrastructure AWS,
        architecture Docker, modèle de données et conception de la
        sécurité \textbf{en profondeur}.
  \item Le \textbf{chapitre 4} détaille la \textit{réalisation et le
        déploiement} : choix technologiques, implémentation,
        construction des images Docker, mise en place de l'instance
        EC2, configuration de Route\,53, du proxy Nginx et des
        certificats Let's Encrypt.
  \item Le \textbf{chapitre 5} expose la phase de \textit{tests
        fonctionnels et d'analyse de sécurité} : scans Nmap, audit
        Nikto, évaluation TLS et synthèse des recommandations.
\end{itemize}

Une \textit{conclusion générale} dresse le bilan du projet et propose
des pistes d'évolution.
